Silicon photonic devices used in photonic integrated circuits (PICs) enable compact optical systems for
communications, sensing, and computing, but practical design still depends on full-wave electromagnetic simulation and
iterative optimization. Finite-difference time domain (FDTD) solvers remain the main tool for accurate prediction
across device geometries and wavelengths~\cite{TafloveHagness2005FDTD, Teixeira2023FDTDMethods, Oskooi2010Meep}, yet
repeated FDTD solves become a bottleneck for parameter sweeps and gradient-based inverse design
loops~\cite{Molesky2018InverseDesign}. This has motivated
machine-learning surrogates for forward prediction and learning-assisted inverse design. A key remaining challenge is
accurate full-field prediction of complex electromagnetic responses: phase must be handled consistently, solutions
should satisfy the wave equation, and models must generalize across device families and wavelengths rather than
overfitting to one geometry class. Moreover, predictions should yield accurate scattering parameters
(S-parameters)---the port-to-port transmission and reflection coefficients that are the primary figures of merit in
photonic device design---without requiring a separate regression head or post-hoc extraction step.

We address this gap with a physics-based complex-valued conditional flow-matching model for fast full-field prediction
of silicon photonic devices. Building on flow matching for generative modeling~\citep{Lipman2022FlowMatching} and
physics-based flow matching for PDE-constrained generation~\citep{Baldan2025PBFM}, we train a complex-valued U-Net to
generate complex field distributions while preserving phase relationships. Physics is enforced through a Helmholtz
residual loss with an interface-aware mask that excludes dielectric boundary pixels where finite-difference stencil
errors dominate, yielding a meaningful compliance metric. S-parameter accuracy is enforced through a dedicated loss
term: at each device port, the predicted field is projected onto the computed fundamental waveguide eigenmode via an
overlap integral, and the resulting complex mode amplitudes are normalized to obtain S-parameters; the loss then
penalizes both magnitude and phase errors relative to FDTD-extracted ground truth, with an amplitude gate that
suppresses supervision on near-zero transmissions. These objectives are introduced through a three-phase training
curriculum---flow matching only, then Helmholtz residual, then S-parameter supervision---that avoids gradient conflicts
between the velocity-matching, PDE-compliance, and port-coupling objectives. We evaluate the approach on canonical PIC
building blocks (bends, tapers, couplers, crossings) and show that a single checkpoint generalizes across device
classes and wavelengths, producing full-field and S-parameter predictions in a single forward pass that would otherwise
require minutes of FDTD simulation.
